% =============================================================================
% APPENDIX BL: REF-EBF: Evidence-Based Framework Guide
% =============================================================================
%
% Category: REF (Reference Materials)
% Module: BCM2_00_META
% Version: 1.0 (January 2026)
% Status: Draft
% Language: English
%
% This appendix explains what EBF is, how it works, and how to use it.
% It serves as the central reference for understanding the framework.
%
% =============================================================================

% =============================================================================
% CROSS-REFERENCE STRUCTURE
% =============================================================================
%
% CHAPTER LINKAGE:
% - Primary Chapter: All chapters (meta-level documentation)
% - Secondary Chapters: Chapter 1 (Introduction), Chapter 9 (Context)
%
% DEPENDENCIES:
% - All CORE appendices (AAA, B, C, V, BBB, AU, AV, AW, HI, IE)
%
% DEPENDENTS:
% - All users of the EBF framework
% - CLAUDE.md (references this appendix)
%
% =============================================================================

\chapter{REF-EBF: Evidence-Based Framework Guide}
\label{app:ref-ebf}

\begin{abstract}
This appendix explains what the Evidence-Based Framework (EBF) is, how it
differs from traditional AI/LLM systems, and how to use it effectively.
EBF is a behavioral-architecture system that models human behavior based on
50+ years of validated research, not a chatbot that generates language.
\end{abstract}

% -----------------------------------------------------------------------------
% QUICK REFERENCE BOX
% -----------------------------------------------------------------------------

\begin{tcolorbox}[colback=blue!5!white, colframe=blue!75!black,
    title=Quick Reference: EBF Framework]

\textbf{Core Question:} How does EBF work and what makes it different?

\textbf{Key Formula:}
\begin{equation}
B = f(N, K, A, W, J)
\end{equation}
where $B$ = Behavior, $N$ = Utility, $K$ = Context, $A$ = Awareness,
$W$ = Willingness, $J$ = Journey.

\textbf{Key Concepts:}
\begin{itemize}[nosep]
\item EBF models \textbf{behavior}, not language
\item \textbf{Smart Data}: 2,300+ curated papers instead of big data
\item \textbf{Context} is the central multiplier for behavior
\item \textbf{8-step EBF Workflow} from question to answer
\end{itemize}

\textbf{Cross-References:} V (CORE-WHEN), AAA (CORE-WHO), C (CORE-WHAT), B (CORE-HOW)
\end{tcolorbox}

% -----------------------------------------------------------------------------
% SCOPE BOX
% -----------------------------------------------------------------------------

\begin{tcolorbox}[
    colback=orange!5!white,
    colframe=orange!75!black,
    title={\textbf{Appendix Scope: Ziel / In-Scope / Out-of-Scope / Constraints}},
    fonttitle=\bfseries
]

\textbf{Ziel (Objective):}
\begin{itemize}[nosep]
    \item Explain what EBF is and how to use it effectively
\end{itemize}

\textbf{In-Scope:}
\begin{itemize}[nosep]
    \item What EBF is (and what it is not)
    \item The scientific basis (papers, theories, cases)
    \item Smart Data vs Big Data philosophy
    \item Context as central multiplier
    \item The EBF Workflow (8 steps)
    \item The 4-layer architecture
\end{itemize}

\textbf{Out-of-Scope (delegiert an andere Appendices):}
\begin{itemize}[nosep]
    \item Mathematical formalization $\rightarrow$ Appendix A (FORMAL-FOUND)
    \item Context dimensions in detail $\rightarrow$ Appendix V (CORE-WHEN)
    \item Utility theory $\rightarrow$ Appendix C (CORE-WHAT)
    \item Parameter values $\rightarrow$ Appendix BBB (CORE-WHERE)
\end{itemize}

\textbf{Lieferobjekte (Deliverables):}
\begin{itemize}[nosep]
    \item Clear understanding of EBF philosophy
    \item Overview of the EBF Workflow
    \item Comparison: EBF vs traditional LLMs
\end{itemize}

\end{tcolorbox}


%%%%%%%%%%%%%%%%%%%%%%%%%%%%%%%%%%%%%%%%%%%%%%%%%%%%%%%%%%%%%%%%%%%%%%%%%%%%%%%
\section{The Fundamental Question}
\label{sec:ref-ebf-fundamental}
%%%%%%%%%%%%%%%%%%%%%%%%%%%%%%%%%%%%%%%%%%%%%%%%%%%%%%%%%%%%%%%%%%%%%%%%%%%%%%%

\begin{tcolorbox}[colback=red!5!white, colframe=red!75!black,
    title=The Fundamental Question]
\textbf{What is EBF and how does it differ from a chatbot?}
\end{tcolorbox}

\subsection{Why This Matters}

Many users approach EBF expecting a chatbot experience---ask a question,
get an answer. But EBF is fundamentally different. Understanding this
difference is crucial for using EBF effectively.

\subsection{The Core Insight}

\begin{tcolorbox}[colback=green!5!white, colframe=green!75!black]
\textbf{EBF is not a chatbot. EBF shows you why people ACT.}

\begin{center}
\begin{tabular}{ll}
ChatGPT tells you & what people \textbf{say}. \\
EBF shows you & why people \textbf{act}. \\
\end{tabular}
\end{center}
\end{tcolorbox}


%%%%%%%%%%%%%%%%%%%%%%%%%%%%%%%%%%%%%%%%%%%%%%%%%%%%%%%%%%%%%%%%%%%%%%%%%%%%%%%
\section{What is EBF?}
\label{sec:ref-ebf-what}
%%%%%%%%%%%%%%%%%%%%%%%%%%%%%%%%%%%%%%%%%%%%%%%%%%%%%%%%%%%%%%%%%%%%%%%%%%%%%%%

EBF is a \textbf{Behavioral-Architecture System} that:

\begin{enumerate}
\item \textbf{Models behavior}, not generates language
\item \textbf{Uses axioms}, not statistics
\item \textbf{Provides causality}, not correlations
\item \textbf{Requires context}, not big data
\item \textbf{Produces explainability}, not plausibility
\end{enumerate}

\subsection{The Mathematical Foundation}

EBF is based on the Behavioral Change Model (BCM 2.0):

\begin{equation}
B = f(N, K, A, W, J)
\end{equation}

where:
\begin{itemize}[nosep]
\item $N$ = Utility structure (individual, collective, identity)
\item $K$ = Context (salience, dynamics)
\item $A$ = Awareness function
\item $W$ = Willingness function
\item $J$ = Behavioral Change Journey
\end{itemize}

\subsection{EBF vs LLM: The Comparison}

\begin{table}[htbp]
\centering
\caption{EBF vs Traditional LLMs}
\label{tab:ebf-vs-llm}
\renewcommand{\arraystretch}{1.3}
\begin{tabular}{@{}lll@{}}
\toprule
\textbf{Aspect} & \textbf{LLM (ChatGPT etc.)} & \textbf{EBF} \\
\midrule
Goal & Text / Answer & Behavior / Impact \\
Basis & Language statistics & Axioms \\
Logic & Correlation & Causality \\
Data & Massive, exogenous & Minimal, contextual \\
Validity & Narrative & Formal \\
Output & Words & Architecture \\
\bottomrule
\end{tabular}
\end{table}


%%%%%%%%%%%%%%%%%%%%%%%%%%%%%%%%%%%%%%%%%%%%%%%%%%%%%%%%%%%%%%%%%%%%%%%%%%%%%%%
\section{The Scientific Basis}
\label{sec:ref-ebf-science}
%%%%%%%%%%%%%%%%%%%%%%%%%%%%%%%%%%%%%%%%%%%%%%%%%%%%%%%%%%%%%%%%%%%%%%%%%%%%%%%

\subsection{50+ Years of Research}

EBF integrates the work of:
\begin{itemize}[nosep]
\item \textbf{Daniel Kahneman} (Nobel Prize 2002) --- System 1 \& 2, Prospect Theory
\item \textbf{Richard Thaler} (Nobel Prize 2017) --- Nudging, Mental Accounting
\item \textbf{Ernst Fehr} --- Social Preferences, Fairness, Trust
\item \textbf{Cass Sunstein} --- Choice Architecture, Libertarian Paternalism
\item \textbf{Dan Ariely} --- Predictable Irrationality
\end{itemize}

\subsection{The Paper Database}

EBF contains \textbf{2,300+ scientific papers} curated by the scientific team
around Prof. Ernst Fehr. These papers are used for:

\begin{enumerate}
\item \textbf{Validating theories} --- 134 theoretical models
\item \textbf{Estimating parameters} --- $\lambda$, $\beta$, $\gamma$, etc.
\item \textbf{Supporting models} --- EBF-specific models like PSF-2.0
\item \textbf{Documenting cases} --- 850+ case studies
\item \textbf{Proving interventions} --- What works where?
\item \textbf{Validating context factors} --- 400+ factors for DACH region
\end{enumerate}


%%%%%%%%%%%%%%%%%%%%%%%%%%%%%%%%%%%%%%%%%%%%%%%%%%%%%%%%%%%%%%%%%%%%%%%%%%%%%%%
\section{Smart Data vs Big Data}
\label{sec:ref-ebf-smartdata}
%%%%%%%%%%%%%%%%%%%%%%%%%%%%%%%%%%%%%%%%%%%%%%%%%%%%%%%%%%%%%%%%%%%%%%%%%%%%%%%

\begin{tcolorbox}[colback=yellow!5!white, colframe=yellow!75!black,
    title=The Core Insight]
\textbf{We don't need millions of your data points. We have 50 years of validated research.}
\end{tcolorbox}

\subsection{Why Papers Are More Valuable Than Raw Data}

\begin{enumerate}
\item \textbf{Validated \& Replicated} --- Each paper is peer-reviewed.
      Meta-analyses aggregate 50-100 studies.
\item \textbf{Causally Identified} --- RCTs, natural experiments, IVs.
      Not just ``X correlates with Y'' but ``X causes Y.''
\item \textbf{Parameters with Uncertainty} --- $\lambda = 2.25$ [1.5, 3.0] with 95\% CI.
\item \textbf{Context Conditions Documented} --- We know WHEN effects apply.
\item \textbf{Mechanisms Explained} --- We understand the WHY, not just the WHAT.
\end{enumerate}

\subsection{What We Need From Clients}

\begin{table}[htbp]
\centering
\caption{What EBF Needs vs Doesn't Need}
\label{tab:data-needs}
\renewcommand{\arraystretch}{1.3}
\begin{tabular}{@{}ll@{}}
\toprule
\textbf{NOT Needed} & \textbf{Needed} \\
\midrule
Millions of transactions & Context (country, industry) \\
Complete customer history & Target group (segment) \\
Personal data & Decision situation \\
Data warehouse access & Desired behavior \\
Months of data preparation & 30 minutes briefing \\
\bottomrule
\end{tabular}
\end{table}

\subsection{Client Data: Helpful, Not Necessary}

\begin{tcolorbox}[colback=blue!5!white, colframe=blue!75!black]
\textbf{Without client data:} Good estimate (literature prior) \\
$\lambda = 2.25$ [1.5, 3.0]

\textbf{With client data:} Better estimate (Bayesian posterior) \\
$\lambda = 2.08$ [1.9, 2.3]

Both are CORRECT. One is more PRECISE.
\end{tcolorbox}


%%%%%%%%%%%%%%%%%%%%%%%%%%%%%%%%%%%%%%%%%%%%%%%%%%%%%%%%%%%%%%%%%%%%%%%%%%%%%%%
\section{Context as Central Multiplier}
\label{sec:ref-ebf-context}
%%%%%%%%%%%%%%%%%%%%%%%%%%%%%%%%%%%%%%%%%%%%%%%%%%%%%%%%%%%%%%%%%%%%%%%%%%%%%%%

\begin{tcolorbox}[colback=red!5!white, colframe=red!75!black,
    title=The Context Principle]
\textbf{Without context $\rightarrow$ Behavior remains potential.} \\
\textbf{With context $\rightarrow$ Behavior becomes real.}
\end{tcolorbox}

\subsection{The Context Formula}

\begin{equation}
N_{eff} = N \times C_{sal} \times C_{compat}
\end{equation}

where:
\begin{itemize}[nosep]
\item $N_{eff}$ = Effective utility
\item $N$ = Base utility
\item $C_{sal}$ = Context salience (how visible is the topic?)
\item $C_{compat}$ = Context compatibility (does utility fit the environment?)
\end{itemize}

\subsection{The 8 Context Dimensions ($\Psi$)}

\begin{table}[htbp]
\centering
\caption{The 8 $\Psi$-Dimensions}
\label{tab:psi-dimensions}
\renewcommand{\arraystretch}{1.3}
\begin{tabular}{@{}cll@{}}
\toprule
\textbf{Symbol} & \textbf{Dimension} & \textbf{Question} \\
\midrule
$\Psi_I$ & Rules & What laws, defaults, regulations apply? \\
$\Psi_S$ & Social & Who is present? What social norms? \\
$\Psi_C$ & Cognitive & Tired? Stressed? Attentive? Motivated? \\
$\Psi_K$ & Culture & What values, traditions, religion? \\
$\Psi_E$ & Resources & How much budget, time, energy? \\
$\Psi_T$ & Time & When? Time pressure? Life phase? \\
$\Psi_M$ & Tools & What technology, infrastructure? \\
$\Psi_F$ & Place & Where physically? Home, office, public? \\
\bottomrule
\end{tabular}
\end{table}

For details, see \textbf{Appendix V (CORE-WHEN)}.


%%%%%%%%%%%%%%%%%%%%%%%%%%%%%%%%%%%%%%%%%%%%%%%%%%%%%%%%%%%%%%%%%%%%%%%%%%%%%%%
\section{The EBF Workflow}
\label{sec:ref-ebf-workflow}
%%%%%%%%%%%%%%%%%%%%%%%%%%%%%%%%%%%%%%%%%%%%%%%%%%%%%%%%%%%%%%%%%%%%%%%%%%%%%%%

Every question goes through the \textbf{EBF Workflow}:

\begin{tcolorbox}[colback=gray!5!white, colframe=gray!75!black,
    title=The EBF Workflow: 8 Steps]

\begin{tabular}{@{}rl@{}}
\textbf{Step 0} & Session starten \\
\textbf{Step 1} & Kontext verstehen \\
\textbf{Step 2} & Modell auswählen \\
\textbf{Step 3} & Parameter bestimmen \\
\textbf{Step 4} & Analyse \& Antwort \\
\textbf{Step 5} & Intervention designen (bei Verhaltenszielen) \\
\textbf{Step 6} & Bericht erstellen \\
\textbf{Step 7} & Ergebnisse sichern \\
\textbf{Step 8} & Qualität prüfen \\
\end{tabular}

\end{tcolorbox}

\subsection{Three Modes}

\begin{table}[htbp]
\centering
\caption{EBF Workflow Modes}
\label{tab:workflow-modes}
\renewcommand{\arraystretch}{1.3}
\begin{tabular}{@{}llll@{}}
\toprule
\textbf{Mode} & \textbf{Duration} & \textbf{Feedback} & \textbf{Output} \\
\midrule
SCHNELL & 10 min & None & $\sim$800 words \\
STANDARD & 45 min & Per step & $\sim$3,000 words \\
TIEF & 2+ hours & Iterative & $\sim$5,000+ words \\
\bottomrule
\end{tabular}
\end{table}


%%%%%%%%%%%%%%%%%%%%%%%%%%%%%%%%%%%%%%%%%%%%%%%%%%%%%%%%%%%%%%%%%%%%%%%%%%%%%%%
\section{The 4-Layer Architecture}
\label{sec:ref-ebf-architecture}
%%%%%%%%%%%%%%%%%%%%%%%%%%%%%%%%%%%%%%%%%%%%%%%%%%%%%%%%%%%%%%%%%%%%%%%%%%%%%%%

\begin{tcolorbox}[colback=purple!5!white, colframe=purple!75!black,
    title=EBF Architecture]

\textbf{Layer 4: User Interface}
\begin{itemize}[nosep]
\item Natural language (``Why don't people save?'')
\item Slash commands (/find-model, /design-intervention)
\end{itemize}

\textbf{Layer 3: LLM Reasoning Engine (Claude)}
\begin{itemize}[nosep]
\item Understands questions in natural language
\item Applies EBF rules (from CLAUDE.md)
\item Executes Python scripts (Code Interpreter)
\item Synthesizes knowledge into answers
\end{itemize}

\textbf{Layer 2: EBF Knowledge Base (Databases)}
\begin{itemize}[nosep]
\item bcm\_master.bib $\rightarrow$ 2,300+ scientific papers
\item theory-catalog.yaml $\rightarrow$ 134 economic models
\item case-registry.yaml $\rightarrow$ 850+ case studies
\item model-registry.yaml $\rightarrow$ EBF-specific models
\item parameter-registry.yaml $\rightarrow$ Validated parameters
\end{itemize}

\textbf{Layer 1: EBF Framework (Rules \& Structure)}
\begin{itemize}[nosep]
\item CLAUDE.md $\rightarrow$ Workflow rules, 10C Framework
\item 10C CORE Questions $\rightarrow$ WHO, WHAT, HOW, WHEN, ...
\item 8 $\Psi$-Dimensions $\rightarrow$ Context analysis
\end{itemize}

\end{tcolorbox}


%%%%%%%%%%%%%%%%%%%%%%%%%%%%%%%%%%%%%%%%%%%%%%%%%%%%%%%%%%%%%%%%%%%%%%%%%%%%%%%
\section{Worked Example}
\label{sec:ref-ebf-example}
%%%%%%%%%%%%%%%%%%%%%%%%%%%%%%%%%%%%%%%%%%%%%%%%%%%%%%%%%%%%%%%%%%%%%%%%%%%%%%%

\subsection{Example: Why Don't People Save for Retirement?}

\paragraph{Question.} ``Why don't people save enough for retirement?''

\paragraph{Step 1: Context.} DACH region, middle-income, working age 35-55.

\paragraph{Step 2: Model.} From theory-catalog: MS-TP-001 (Hyperbolic Discounting),
MS-RD-003 (Exponential Growth Bias).

\paragraph{Step 3: Parameters.} From literature:
\begin{itemize}[nosep]
\item $\beta = 0.70$ (present bias)
\item $\delta = 0.95$ (long-run discount factor)
\end{itemize}

\paragraph{Step 4: Analysis.} Present bias ($\beta < 1$) leads to
procrastination. Exponential growth bias underestimates compound interest.

\paragraph{Step 5: Intervention.} Opt-out default enrollment
(see case-registry: CASE-002, Madrian \& Shea 2001).

\paragraph{Result.} With opt-out defaults, participation increases from
$\sim$50\% to $>$90\%.


%%%%%%%%%%%%%%%%%%%%%%%%%%%%%%%%%%%%%%%%%%%%%%%%%%%%%%%%%%%%%%%%%%%%%%%%%%%%%%%
\section{Implications for Practice}
\label{sec:ref-ebf-practice}
%%%%%%%%%%%%%%%%%%%%%%%%%%%%%%%%%%%%%%%%%%%%%%%%%%%%%%%%%%%%%%%%%%%%%%%%%%%%%%%

\begin{tcolorbox}[colback=green!5!white, colframe=green!75!black,
    title=Practical Implications]
\begin{enumerate}
\item \textbf{Context First:} Always define the context before analyzing behavior.
      The same person acts differently in different contexts.
\item \textbf{Use the Workflow:} Follow the 8-step EBF Workflow for every
      substantive question to ensure consistency and completeness.
\item \textbf{Trust the Parameters:} The 2,300+ papers provide validated
      parameter values. You don't need to re-estimate from scratch.
\item \textbf{Client Data Optional:} Client data improves precision but is
      not required. You can start immediately with literature priors.
\end{enumerate}
\end{tcolorbox}


%%%%%%%%%%%%%%%%%%%%%%%%%%%%%%%%%%%%%%%%%%%%%%%%%%%%%%%%%%%%%%%%%%%%%%%%%%%%%%%
\section{How EBF Learns}
\label{sec:ref-ebf-learning}
%%%%%%%%%%%%%%%%%%%%%%%%%%%%%%%%%%%%%%%%%%%%%%%%%%%%%%%%%%%%%%%%%%%%%%%%%%%%%%%

\begin{tcolorbox}[colback=yellow!5!white, colframe=yellow!75!black,
    title=The Learning Principle]
\textbf{EBF learns from validated evidence, not from unstructured data.}
\end{tcolorbox}

\subsection{The 5-Database Learning Loop}

EBF has a continuous learning loop built into its architecture:

\begin{enumerate}
\item \textbf{Session Database} --- Every analysis session is saved with
      question, model, parameters, and results.
\item \textbf{Model Registry} --- New models are registered and validated
      against literature.
\item \textbf{Intervention Registry} --- Projects track predictions vs.
      actual results.
\item \textbf{Parameter Registry} --- Parameters are updated based on
      real-world outcomes (Bayesian updating).
\item \textbf{Case Registry} --- Successful interventions become new cases
      for future reference.
\end{enumerate}

\begin{tcolorbox}[colback=gray!5!white, colframe=gray!75!black]
\begin{center}
\textbf{The Learning Cycle:}

New Question $\rightarrow$ Find Similar Cases $\rightarrow$ Design Model
$\rightarrow$ Make Prediction $\rightarrow$ Execute $\rightarrow$
Measure Results $\rightarrow$ Update Parameters $\rightarrow$ Add New Case
\end{center}
\end{tcolorbox}

\subsection{Bayesian Updating}

When client data becomes available, EBF updates its parameters:

\begin{equation}
\theta_{posterior} = \frac{P(Data | \theta) \cdot P(\theta)}{P(Data)}
\end{equation}

\begin{itemize}[nosep]
\item \textbf{Prior:} Literature estimate (e.g., $\lambda = 2.25$ [1.5, 3.0])
\item \textbf{Likelihood:} Client data (e.g., observed loss aversion)
\item \textbf{Posterior:} Updated estimate (e.g., $\lambda = 2.08$ [1.9, 2.3])
\end{itemize}

\subsection{Knowledge Growth}

The EBF knowledge base grows through:

\begin{table}[htbp]
\centering
\caption{How EBF Knowledge Grows}
\label{tab:knowledge-growth}
\renewcommand{\arraystretch}{1.3}
\begin{tabular}{@{}lll@{}}
\toprule
\textbf{Source} & \textbf{What is learned} & \textbf{Where stored} \\
\midrule
New papers & New theories, parameters & bcm\_master.bib \\
Client projects & Real-world outcomes & intervention-registry \\
Session feedback & Model refinements & model-registry \\
Case studies & What works where & case-registry \\
User corrections & Error patterns & lessons\_learned.md \\
\bottomrule
\end{tabular}
\end{table}

\subsection{What EBF Does NOT Learn From}

\begin{tcolorbox}[colback=red!5!white, colframe=red!75!black]
\textbf{EBF does NOT learn from:}
\begin{itemize}[nosep]
\item Unvalidated user opinions
\item Correlation-only data
\item Single anecdotes without replication
\item Data without context documentation
\end{itemize}

\textbf{Why?} To maintain scientific integrity. Every piece of knowledge
must be traceable to validated research or controlled experiments.
\end{tcolorbox}


%%%%%%%%%%%%%%%%%%%%%%%%%%%%%%%%%%%%%%%%%%%%%%%%%%%%%%%%%%%%%%%%%%%%%%%%%%%%%%%
\section{Open Issues and Future Directions}
\label{sec:ref-ebf-open}
%%%%%%%%%%%%%%%%%%%%%%%%%%%%%%%%%%%%%%%%%%%%%%%%%%%%%%%%%%%%%%%%%%%%%%%%%%%%%%%

\begin{enumerate}
\item \textbf{Cross-cultural validation:} Most parameters are validated
      in WEIRD (Western, Educated, Industrialized, Rich, Democratic) contexts.
      Expansion to other cultural contexts is ongoing.

\item \textbf{Dynamic parameters:} Some parameters (e.g., trust levels)
      change over time. Better temporal modeling is needed.

\item \textbf{Interaction effects:} The interaction between context
      dimensions ($\Psi_i \times \Psi_j$) needs more empirical validation.

\item \textbf{Individual differences:} How to incorporate personality
      traits and individual heterogeneity more systematically.
\end{enumerate}


%%%%%%%%%%%%%%%%%%%%%%%%%%%%%%%%%%%%%%%%%%%%%%%%%%%%%%%%%%%%%%%%%%%%%%%%%%%%%%%
\section{Glossary of Symbols}
\label{sec:ref-ebf-glossary}
%%%%%%%%%%%%%%%%%%%%%%%%%%%%%%%%%%%%%%%%%%%%%%%%%%%%%%%%%%%%%%%%%%%%%%%%%%%%%%%

\begin{table}[htbp]
\centering
\caption{Key Symbols in EBF}
\label{tab:symbols}
\renewcommand{\arraystretch}{1.3}
\begin{tabular}{@{}cll@{}}
\toprule
\textbf{Symbol} & \textbf{Name} & \textbf{Description} \\
\midrule
$B$ & Behavior & Observable action or decision \\
$N$ & Utility & Net utility structure (individual + social + identity) \\
$K$ & Context & Situational factors ($\Psi$-dimensions) \\
$A$ & Awareness & Knowledge and attention function \\
$W$ & Willingness & Readiness to act \\
$J$ & Journey & Position in Behavioral Change Journey \\
$\Psi$ & Psi & Context dimensions (8 total) \\
$\lambda$ & Lambda & Loss aversion coefficient (typically $\sim$2.25) \\
$\beta$ & Beta & Present bias (hyperbolic discounting) \\
$\gamma$ & Gamma & Complementarity coefficient \\
$C_{sal}$ & Salience & Context salience factor \\
$C_{compat}$ & Compatibility & Context compatibility factor \\
$N_{eff}$ & Effective utility & $N \times C_{sal} \times C_{compat}$ \\
\bottomrule
\end{tabular}
\end{table}

For the complete glossary, see \textbf{Appendix G (Glossary)}.


%%%%%%%%%%%%%%%%%%%%%%%%%%%%%%%%%%%%%%%%%%%%%%%%%%%%%%%%%%%%%%%%%%%%%%%%%%%%%%%
\section{Critical Foundations}
\label{sec:ref-ebf-foundations}
%%%%%%%%%%%%%%%%%%%%%%%%%%%%%%%%%%%%%%%%%%%%%%%%%%%%%%%%%%%%%%%%%%%%%%%%%%%%%%%

The EBF learning mechanisms are grounded in established scientific literature:

\subsection{Bayesian Updating}

The Bayesian approach to parameter updating is foundational:
\begin{itemize}[nosep]
\item \textbf{Gelman et al. (2017):} ``The Prior Can Often Only Be Understood
      in the Context of the Likelihood'' --- establishes that priors must be
      interpretable within the model context \citep{gelman2017prior}.
\item \textbf{Stefan et al. (2025):} ``Using Large Language Models to Suggest
      Informative Prior Distributions'' --- validates LLM-based prior
      elicitation \citep{scientificreports2025llmpriors}.
\end{itemize}

\subsection{Field Experiments and Validation}

Real-world validation through field experiments:
\begin{itemize}[nosep]
\item \textbf{Banerjee \& Duflo (2017):} \textit{Handbook of Field Experiments}
      --- comprehensive methodology for field validation \citep{banerjee2017handbook}.
\item \textbf{List (2020):} ``Field Experiments for Policy Analysis'' ---
      framework for translating lab findings to field \citep{list2020field}.
\item \textbf{Hallsworth et al. (2017):} ``The Behavioralist as Tax Collector''
      --- landmark field experiment validating behavioral interventions
      \citep{hallsworth2017}.
\end{itemize}

\subsection{Parameter Estimation}

Foundational work on behavioral parameters:
\begin{itemize}[nosep]
\item \textbf{Kahneman \& Tversky (1991):} ``Loss Aversion in Riskless Choice''
      --- establishes $\lambda \approx 2.25$ \citep{kahneman1991loss}.
\item \textbf{Gächter \& Klinowski (2022):} ``Individual-Level Loss Aversion''
      --- heterogeneity in loss aversion parameters \citep{PAP-gaechter2022lossaversion}.
\item \textbf{Benartzi \& Thaler (1995):} ``Myopic Loss Aversion'' ---
      context-dependent loss aversion \citep{benartzi1995myopic}.
\end{itemize}

\subsection{Machine Learning Integration}

Econometric foundations for ML in behavioral economics:
\begin{itemize}[nosep]
\item \textbf{Athey \& Imbens (2019):} ``Machine Learning Methods That
      Economists Should Know About'' \citep{athey2019machine}.
\item \textbf{Mullainathan \& Spiess (2017):} ``Machine Learning: An Applied
      Econometric Approach'' \citep{mullainathan2017machine}.
\end{itemize}

\subsection{Learning from Interventions}

Evidence that interventions produce measurable, learnable outcomes:
\begin{itemize}[nosep]
\item \textbf{Hussam et al. (2022):} ``Rational Habit Formation: Experimental
      Evidence from Handwashing in India'' --- demonstrates measurable
      behavioral persistence \citep{HussamEtAl2022}.
\item \textbf{Jones, Molitor \& Reif (2025):} ``How Do Habits Form?'' ---
      habit formation dynamics from field experiments \citep{JonesMolitorReif2025}.
\end{itemize}


%%%%%%%%%%%%%%%%%%%%%%%%%%%%%%%%%%%%%%%%%%%%%%%%%%%%%%%%%%%%%%%%%%%%%%%%%%%%%%%
\section{Summary}
\label{sec:ref-ebf-summary}
%%%%%%%%%%%%%%%%%%%%%%%%%%%%%%%%%%%%%%%%%%%%%%%%%%%%%%%%%%%%%%%%%%%%%%%%%%%%%%%

\begin{tcolorbox}[colback=green!10!white, colframe=green!75!black,
    title=\textbf{Appendix BL: Summary}]

\textbf{Core Message:} EBF is not a chatbot---it's a behavioral-architecture
system that models WHY people act.

\textbf{Key Principles:}
\begin{itemize}[nosep]
\item Smart Data: 2,300+ papers instead of big data
\item Context is the central multiplier
\item 8-step EBF Workflow
\item 4-layer architecture
\end{itemize}

\textbf{The Motto:}
\begin{quote}
``Scientific coherence meets cognitive scalability.''
\end{quote}

\end{tcolorbox}


%%%%%%%%%%%%%%%%%%%%%%%%%%%%%%%%%%%%%%%%%%%%%%%%%%%%%%%%%%%%%%%%%%%%%%%%%%%%%%%
\section{References}
\label{sec:ref-ebf-references}
%%%%%%%%%%%%%%%%%%%%%%%%%%%%%%%%%%%%%%%%%%%%%%%%%%%%%%%%%%%%%%%%%%%%%%%%%%%%%%%

\nocite{bcm_master}

\begin{tcolorbox}[colback=blue!5!white, colframe=blue!75!black]
\textbf{Master Bibliography:} See \texttt{bcm\_master.bib} for complete references.
\end{tcolorbox}

\vspace{1em}
\hrule
\vspace{0.5em}

\noindent\textit{Cross-references:}
\begin{itemize}[nosep]
\item Appendix G (Glossary) --- Master glossary of all EBF terms
\item Appendix V (CORE-WHEN) --- Context dimensions in detail
\item Appendix AAA (CORE-WHO) --- The Welfare Hierarchy
\item Appendix C (CORE-WHAT) --- Utility theory
\item Appendix B (CORE-HOW) --- Complementarity
\item Appendix BBB (CORE-WHERE) --- Parameter repository
\item Appendix HHH (METHOD-TOOLKIT) --- Intervention design
\end{itemize}

\noindent\textit{Master Bibliography:} \texttt{bcm\_master.bib}

% =============================================================================
% END OF APPENDIX BL
% =============================================================================
